\documentclass[12pt,a4paper]{article}

\usepackage[utf8]{inputenc}
\usepackage[italian]{babel}
\usepackage{Alegreya}
\usepackage[T1]{fontenc}
\usepackage{geometry}
\usepackage{setspace}
\usepackage{titlesec}
\usepackage{enumitem}

\geometry{margin=2.5cm}
\onehalfspacing

\titleformat{\section}{\Large\bfseries}{\thesection.}{0.5em}{}
\titleformat{\subsection}{\large\bfseries}{\thesubsection.}{0.5em}{}
\titleformat{\subsubsection}{\normalsize\bfseries}{\thesubsubsection.}{0.5em}{}

\setlist[itemize]{noitemsep, topsep=0pt}

\title{\Huge\bfseries LAZZARO IN RIPRESA DI VENTI\\
\Large Perché siamo qui?\\
\normalsize Utopia-Anacronismo: di quei fasci di luce nelle tenebre irradiati dal \emph{Maestro delle Ombre}}
\author{}
\date{}

\begin{document}

\maketitle

\begin{abstract}
Appunti sparsi nella collezione di \emph{historìa} riguardo Mario Bertoncini nel condurre l'apertura dell'introduzione all'ascolto per LAZZARO IN RIPRESA DI VENTI
\end{abstract}

\tableofcontents
\newpage

\section{Ragionamenti in forma di dialogo - p. 212}

Molto s’è detto e scritto, in tutti i tempi ma soprattutto nel nostro, dei limiti,
degli scopi, dei destini e persino della morte dell’arte e senza dubbio, a dispetto di
quanto detto o scritto con maggiore o minore autorevolezza, d’arte si continua a
parlare e spesso a torto o a ragione, soprattutto a farne.
Non sta a noi, che d’arte e per l’arte viviamo, decidere quando essa sia legitti-
ma, quando cioè, se ci riferiamo a epoche forse più felici o meno travagliate della
nostra, essa appaia in sintonia con le esigenze della società circostante: anche se
esattamente questo sembrerebbe mancare oggi in un mondo oppresso da interrogati-
vi irrisolti, da tragedie ricorrenti e dalle angosce d’un quotidiano privo di adeguata
soluzione.\\

[...]\\

Quale può dunque essere l’alternativa desiderabile per un’arte amata oppure
odiata ma comunque a dispetto di tutto sempre presente tuttavia in un’epoca sia
pure distratta e turbolenta come quella nostra? O meglio, per scoprire le carte, qual
è la soluzione che noi caldeggiamo, volendo portare il suono del nostro immaginario
a confronto con uno spazio più vasto, come è quello percorso dal vento o dalla magia
delle campane?\\


\subsection{VI. Objet trouvé. Alcune osservazioni circa l’uso o la lettura delle sculture di suono - p.267}

Nel tipo particolare di musica che si vuole descrivere qui, suonare, ese-
guire, è un’azione idealmente e metaforicamente assimilabile a quella
di scavo. Come l’archeologo libera accuratamente dalle incrostazioni
e dai detriti che lo ricoprono il reperto ambito, cercato ma non cono-
sciuto in anticipo, così all’esecutore attento si chiede in questo caso
di scegliere fra i tanti percorsi possibili lungo la densa compagine
di corde quell’unico gesto o quella serie di gesti già esplicitamente
catalogati nelle note per l’esecuzione.
\\
[...]\\

L’insieme delle misure necessarie alla
realizzazione, allo sviluppo sonoro dell’oggetto, sono indispensabili per
l’esecutore che voglia tradurre in atto una interpretazione adeguata di esso,
ma non sono da confondere con un piano univoco di sviluppo temporale.
E qui attenzione: tutte le volte che io stesso definisco graficamente (o
verbalmente) una determinata partitura, non detto in modo univoco l’unica
interpretazione dell’oggetto, ma soltanto una delle molteplici interpretazioni
di esso e sono quindi uno dei tanti possibili interpreti della mia stessa
composizione.\\


Quale potrà essere quindi la definitiva sistemazione dei rispettivi ruoli
del compositore, dell’interprete e del fruitore e quale la legittima
forma con cui trasmettere stabilmente e validamente un pensiero
musicale?\\

[...]\\

Ancora un cenno sulla lettura dell’oggetto sonoro. Abbiamo af-
fermato che l’osservatore–ricettore–esecutore dell’oggetto, nell’atto
cosciente dell’applicare alla scultura di suono una successione qual-
siasi dei modelli esecutivi memorizzati, nel rivelarlo scopre incon-
sciamente un ductus sonoro agente sulla propria immaginazione; e la
sorpresa che ne deriva agisce su quella con la stessa vivida sensazione
che accompagna il sogno. E come nel sogno, anche nel processo
eolico di lettura, in quel groviglio di corde delle slegate immagini
oniriche si ricollegano continuamente ad una rete di ricordi allusivi la
cui assurdità fantastica viene proposta in modo momentaneamente
credibile. È possibile — o probabile — che fra tali immagini sonore
affiori un subitaneo accenno ad un inciso noto, seppure confuso e
sfocato da un salutare tessuto magmatico contrastante: tale inciso, a
nostro parere, è l’objet trouvé, appunto.

\subsubsection{Una nuova Utopia: Vele per Valeska - Dialogo Duodecimo p289}

Gabbia de’ matti è ’l mondo. . .\\
— Tommaso Campanella\\

MAESTRO: Dove eravamo rimasti?. . .\\
BREMONTE: Lei stava parlando dello svolgimento di Vele per Valeska
quando la conversazione fu dirottata verso il poststrutturalismo e la
decostruzione di Derrida. . .\\
DR. WIESEL: Per favore, non ricominciamo! Mi aiuti, invece, Bremon-
te: desidererei che il Maestro esponesse nei particolari quel lavoro
— lavoro o sogno che sia —, ignorando la «gabbia de’ matti» cui si
riferiva Tommaso Campanella. . .\\
MAESTRO: Beh, ero sul punto di dire che quel progetto campato in aria,
anziché ricalcare utopicamente l’ambiziosa idea originale delle tre grandi
arpe eolie da installare in montagna, all’aperto — un’idea che dati i costi
sono riuscito poi a eseguire soltanto in forma ridotta —, nella leggerezza
del sogno si presentò in una veste ancora più impegnativa e dispendiosa.
\emph{Dato che per me quell’idea era irrealizzabile, perché limitarne la portata
per amore d’un banale realismo?}\\

\subsubsection{Emanuele Severino - La filosofia Futura}

L'interpretazione è la volontà che il mondo abbia un senso.\\
Hérmes, il messaggero. \\
Unifica, attiva, le due dimensioni: \emph{reale} e \emph{ideale}.\\

\section{UTOPICI}

\subsection{LAZZARO}

\subsubsection{MARIO}

\subsubsection{Ingeborg Bachmann: Letteratura come Utopia}


\end{document}
